\textit{\textbf{Important:}
PyTV is intended for research, education, and rapid prototyping workflows.
Production deployment requires additional validation and process controls.}

\subsection{Python Execution Trust Model}
PyTV executes generated Python when \texttt{--run-py} or \texttt{--run-py-del} is used.
Any Python embedded in \texttt{.pytv} or preamble files runs with the invoking user's permissions.
Treat templates and preambles as trusted code.

\subsection{Instantiation Scope}
\texttt{<INST>} blocks support single-level instantiation extraction per module conversion.
Nested \texttt{<INST>} blocks are rejected.
Hierarchical assembly across modules should be implemented by a higher-level orchestrator using emitted \texttt{.inst} files.

\subsection{Verilog Parsing Model}
PyTV is line-oriented and template-oriented.
It does not perform full HDL semantic checks.
Backtick substitution and emitted text correctness remain the user's responsibility.

\subsection{Feature-Specific Behavior}
\begin{itemize}
  \item With feature \texttt{inst} disabled, instantiation parsing/output is unavailable.
  \item The optional \texttt{parameters} field in an \texttt{<INST>} block is stored in \texttt{.inst}
  metadata but is not rendered into Verilog module instantiation text.
\end{itemize}

\subsection{License}
PyTV is distributed under GPL-3.0-or-later.
Downstream usage must comply with the GPL terms.
